% GENERATED FILE. Edit canonical lessons, then run npm run generate:books.
% canonical-source-hash: fnv1a64:f781ab7475228ecd
% canonical-lessons: RU-C449-parikmakherskaya, RU-C449-strizhka, RU-C449-strich, RU-C449-britsya, RU-C449-chyolka

\chapter{A haircut}
\label{ch:ru-a-haircut}
\hlchaptermodality{449}

\begin{chapteropening}
\textbf{By the end of this chapter:} \emph{I can get a haircut: the hairdresser's, a haircut, cutting hair, shaving and a fringe.}

Complete the last lesson of chapter 449: I can get a haircut: the hairdresser's, a haircut, cutting hair, shaving and a fringe.
\end{chapteropening}

\section[\texorpdfstring{\ru{парикмахерская} (parikmákherskaya)}{парикмахерская (parikmákherskaya)}]{\ru{парикмахерская} (parikmákherskaya) — a hairdresser's, a barber's}

\label{lesson:RU-C449-parikmakherskaya}



\begin{quote}
Before the new one: say the Russian for to sew up, then the Russian for to sew on.
\end{quote}

\subsection*{You'll want to know: \ru{парикмахерская}}
\textbf{\ru{парикмахерская}} (\emph{parikmákherskaya}) — \textquotedblleft{}a hairdresser's, a barber's\textquotedblright{}.

\begin{quote}
\ru{Парикмахерская} \ru{работает} \ru{до} \ru{восьми}.
\end{quote}

\begin{quote}
\emph{Parikmákherskaya rabótayet do vosmí.}
\end{quote}

\begin{quote}
The hairdresser's is open until eight.
\end{quote}

An adjective used as a noun.

\begin{cousinweb}
\emph{\ru{парикмахерская}} is built on \emph{\ru{парикмахер}} 'a hairdresser' with an adjective ending; \emph{\ru{парикмахер}} was borrowed from German \emph{Perückenmacher} 'a wig-maker' (EtymDB).
\end{cousinweb}

\begin{grammarlens}[title={What you've built}]
One more word for telling, explaining and describing at B1.
\end{grammarlens}

\subsection*{Your turn}
\begin{itemize}
\raggedright
  \item \emph{Say it:} \emph{parikmákherskaya}
  \item \emph{Say it:} \emph{parikmákherskaya}, once more
  \item \emph{Say it:} a sentence of your own with \emph{parikmákherskaya}
  \item \emph{Recall:} say the Russian for to sew up, then the Russian for to sew on, then say \emph{parikmákherskaya} again
\end{itemize}

\begin{tcolorbox}[breakable,colback=teal!4,colframe=teal!35!black,title={Before you move on}]
What does \ru{парикмахерская} mean? (\textquotedblleft{}A hairdresser's, a barber's\textquotedblright{}.) Say it once more.
\end{tcolorbox}
\section[\texorpdfstring{\ru{стрижка} (strízhka)}{стрижка (strízhka)}]{\ru{стрижка} (strízhka) — a haircut}

\label{lesson:RU-C449-strizhka}



\begin{quote}
Before the new one: say the Russian for to sew on, then the Russian for a hairdresser's.
\end{quote}

\subsection*{You'll want to know: \ru{стрижка}}
\textbf{\ru{стрижка}} (\emph{strízhka}) — \textquotedblleft{}a haircut\textquotedblright{}.

\begin{quote}
\ru{Тебе} \ru{идёт} \ru{новая} \ru{стрижка}.
\end{quote}

\begin{quote}
\emph{Tebé idyót nóvaya strízhka.}
\end{quote}

\begin{quote}
The new haircut suits you.
\end{quote}

Feminine.

\begin{cousinweb}
\emph{\ru{стрижка}} is built on \emph{\ru{стричь}} 'to cut hair, to shear' with \emph{-\ru{ка}}. \emph{\ru{Стричь}} is inherited from Proto-Slavic \emph{striťi} (EtymDB).
\end{cousinweb}

\begin{grammarlens}[title={What you've built}]
One more word for telling, explaining and describing at B1.
\end{grammarlens}

\subsection*{Your turn}
\begin{itemize}
\raggedright
  \item \emph{Say it:} \emph{strízhka}
  \item \emph{Say it:} \emph{strízhka}, once more
  \item \emph{Say it:} a sentence of your own with \emph{strízhka}
  \item \emph{Recall:} say the Russian for to sew on, then the Russian for a hairdresser's, then say \emph{strízhka} again
\end{itemize}

\begin{tcolorbox}[breakable,colback=teal!4,colframe=teal!35!black,title={Before you move on}]
What does \ru{стрижка} mean? (\textquotedblleft{}A haircut\textquotedblright{}.) Say it once more.
\end{tcolorbox}
\section[\texorpdfstring{\ru{стричь} (strich)}{стричь (strich)}]{\ru{стричь} (strich) — to cut (hair), to trim}

\label{lesson:RU-C449-strich}



\begin{quote}
Before the new one: say the Russian for a hairdresser's, then the Russian for a haircut.
\end{quote}

\subsection*{You'll want to know: \ru{стричь}}
\textbf{\ru{стричь}} (\emph{strich}) — \textquotedblleft{}to cut (hair), to trim\textquotedblright{}.

\begin{quote}
\ru{Мастер} \ru{стрижёт} \ru{меня} \ru{уже} \ru{пять} \ru{лет}.
\end{quote}

\begin{quote}
\emph{Máster strizhyót menyá uzhé pyat let.}
\end{quote}

\begin{quote}
The hairdresser has been cutting my hair for five years now.
\end{quote}

Imperfective; the perfective is \emph{\ru{постричь}}. Present \emph{\ru{стригу}, \ru{стрижёшь}, \ru{стригут}}; past \emph{\ru{стриг}}. \emph{\ru{Стричься}} — to have one's hair cut.

\begin{cousinweb}
\emph{\ru{стричь}} is inherited from Proto-Slavic \emph{striťi} (EtymDB), which EtymDB takes back to a PIE root \emph{streyg-}.
\end{cousinweb}

\begin{grammarlens}[title={What you've built}]
One more word for telling, explaining and describing at B1.
\end{grammarlens}

\subsection*{Your turn}
\begin{itemize}
\raggedright
  \item \emph{Say it:} \emph{strich}
  \item \emph{Say it:} \emph{strich}, once more
  \item \emph{Say it:} a sentence of your own with \emph{strich}
  \item \emph{Recall:} say the Russian for a hairdresser's, then the Russian for a haircut, then say \emph{strich} again
\end{itemize}

\begin{tcolorbox}[breakable,colback=teal!4,colframe=teal!35!black,title={Before you move on}]
What does \ru{стричь} mean? (\textquotedblleft{}To cut (hair), to trim\textquotedblright{}.) Say it once more.
\end{tcolorbox}
\section[\texorpdfstring{\ru{бриться} (brítsya)}{бриться (brítsya)}]{\ru{бриться} (brítsya) — to shave}

\label{lesson:RU-C449-britsya}



\begin{quote}
Before the new one: say the Russian for a haircut, then the Russian for to cut.
\end{quote}

\subsection*{You'll want to know: \ru{бриться}}
\textbf{\ru{бриться}} (\emph{brítsya}) — \textquotedblleft{}to shave\textquotedblright{}.

\begin{quote}
\ru{Отец} \ru{бреется} \ru{каждое} \ru{утро}.
\end{quote}

\begin{quote}
\emph{Otéts bréyetsya kázhdoye útro.}
\end{quote}

\begin{quote}
Father shaves every morning.
\end{quote}

Imperfective; perfective \emph{\ru{побриться}}. Present \emph{\ru{бреюсь}, \ru{бреешься}}.

\begin{cousinweb}
\emph{\ru{бриться}} is \emph{\ru{брить}} 'to shave' with \emph{-\ru{ся}} 'oneself'. \emph{\ru{Брить}} is inherited from Proto-Slavic \emph{briti} (EtymDB).
\end{cousinweb}

\begin{grammarlens}[title={What you've built}]
One more word for telling, explaining and describing at B1.
\end{grammarlens}

\subsection*{Your turn}
\begin{itemize}
\raggedright
  \item \emph{Say it:} \emph{brítsya}
  \item \emph{Say it:} \emph{brítsya}, once more
  \item \emph{Say it:} a sentence of your own with \emph{brítsya}
  \item \emph{Recall:} say the Russian for a haircut, then the Russian for to cut, then say \emph{brítsya} again
\end{itemize}

\begin{tcolorbox}[breakable,colback=teal!4,colframe=teal!35!black,title={Before you move on}]
What does \ru{бриться} mean? (\textquotedblleft{}To shave\textquotedblright{}.) Say it once more.
\end{tcolorbox}
\section[\texorpdfstring{\ru{чёлка} (chyólka)}{чёлка (chyólka)}]{\ru{чёлка} (chyólka) — a fringe (of hair), bangs}

\label{lesson:RU-C449-chyolka}



\begin{quote}
Before the new one: say the Russian for to cut, then the Russian for to shave.
\end{quote}

\subsection*{You'll want to know: \ru{чёлка}}
\textbf{\ru{чёлка}} (\emph{chyólka}) — \textquotedblleft{}a fringe (of hair), bangs\textquotedblright{}.

\begin{quote}
\ru{У} \ru{девочки} \ru{длинная} \ru{чёлка}.
\end{quote}

\begin{quote}
\emph{U dévochki dlínnaya chyólka.}
\end{quote}

\begin{quote}
The girl has a long fringe.
\end{quote}

Feminine. In American English, bangs.

\begin{cousinweb}
\emph{\ru{чёлка}} is usually derived from \emph{\ru{чело}} 'a brow', which is inherited from Proto-Slavic \emph{čelo} 'a forehead' (EtymDB): the hair that falls on the brow. EtymDB does not list \emph{\ru{чёлка}} itself.
\end{cousinweb}

\begin{grammarlens}[title={What you've built}]
One more word for telling, explaining and describing at B1.
\end{grammarlens}

\subsection*{Your turn}
\begin{itemize}
\raggedright
  \item \emph{Say it:} \emph{chyólka}
  \item \emph{Say it:} \emph{chyólka}, once more
  \item \emph{Say it:} a sentence of your own with \emph{chyólka}
  \item \emph{Recall:} say the Russian for to cut, then the Russian for to shave, then say \emph{chyólka} again
\end{itemize}

\begin{tcolorbox}[breakable,colback=teal!4,colframe=teal!35!black,title={Before you move on}]
What does \ru{чёлка} mean? (\textquotedblleft{}A fringe (of hair), bangs\textquotedblright{}.) Say it once more.
\end{tcolorbox}
